\section{Why deleting atomic case introduction is not an equivalence}

Consider the native Eq/$D$ calculus with closure clauses $Q$. Keep
arithmetic computation, equality transport, paired elimination, closure
specialization, and identity/availability. Modify paired introduction so
that its local assumption is always $E(t,z)$ when the matrix is a literal
equation, and always $B(t,0)$ otherwise. This removes the independent
case/progress choice for atomic matrices.

There are two sound interpretations on standard natural-number arithmetic.
In the first, every $D$ has its usual closure reading
\[
  D_{x,y}B=\exists x\bigl(B(x,0)\to\forall y\,B(x,y)\bigr).
\]
In the second, interpret literal equations by actual equality, interpret
$Q_yB$ by $B(0)\to\forall yB(y)$, and define $D$ recursively by
\[
 (D_{x,y}B)^\dagger=
 \begin{cases}
 \exists x\,\forall y\bigl(B(x,y)\to B(x,Sy)\bigr),
   &B\text{ a literal equation},\\
 \exists x\bigl(B^\dagger(x,0)\to\forall yB^\dagger(x,y)\bigr),
   &B\text{ compound}.
 \end{cases}
\]
Both interpretations validate every remaining rule. In particular, an
atomic progress witness supplies its closure clause by ordinary induction
on the standard natural numbers, so paired elimination is sound.
Compound case introduction is sound because every positive natural number
is a successor. Equality transport preserves the literal/compound
classification and respects the displayed interpretations.

Let $E(y)$ be $y=1$, with an unused existential variable $x$. In the first
interpretation, $D_{x,y}E(y)$ is true because $E(0)$ is false. In the second
interpretation it is false, because $E(1)$ is true and $E(2)$ is false.
The original calculus proves this formula by atomic case introduction:
under its base assumption $0=1$, explosion supplies the required successor
equation. Hence the restricted calculus is not equivalent to the original
native calculus.

Padding an equation into a compound formula may recover an ordinary
logical interface, but this does not repair the missing native round trip.
If atomic $D$ is decoded as closure, its closure formula cannot be
translated back to the native $D$ in the second interpretation. If it is
decoded as progress, the native $D$ cannot be translated forward to its
progress formula in the first interpretation. Matrix congruence with an
atomic target also cannot be assumed: its usual proof uses precisely the
atomic case introduction being removed.

This is a limitation of the proposed deletion with the displayed natural
interpretations, not a theorem excluding every possible unrelated
bi-interpretation of every redesign. A possible new design could treat
atomic $D$ certificates as auxiliary judgments only and restrict genuine
object formulas to a padded language. Such a design must specify and
charge its object-language validation, rather than silently counting all
accepted raw formulas as object theorems.
